\scanpage{361}\mspage{277}
En d'autres termes, on peut dire que
le foncteur d'induction canonique
\[
 C\unsure{-|\longrightarrow}\mathrm{Ind}_{\pi_0}C=M
\]
est \uline{2-universel} parmi les foncteurs de
$C$ dans des catégories $N$ satisfaisant $L_{\pi_0}$,
celles-ci \deleted{\unsure{flèches}} étant considérées comme
$0$-objets d'une $2$-catégorie où les
$1$-\deleted{objets}\unsure{flèches} sont les foncteurs $N\longrightarrow N'$
$\pi_0$-accessibles, i.e. commu\deleted{\unsure{tent}}\unsure{tants} aux
$\unsure{\underline{\lim}|\underrightarrow{\lim}}$ \unsure{gdes|grdes|grandes} devant $\pi_0$.

Cet énoncé est l'analogue, \inserted{pour $C$ \inserted{petite} quelconque,} de 4.4.16
(page \unsure{101}), valable pour $C$ stable par
$\underline{\lim}$ de cardinal $\leq\pi_0$. Il y a bien
sûr d'autres variantes, en relation avec
la \deleted{stabilité} commutation à d'autres
types de $\underline{\lim}$ ($\underline{\lim}$ filtrantes, sommes,
limites inductives finies etc) envisagés
dans la section 4.11. précédente.

La dém. de 4.12.3 est \unsure{triviale}.
% archive sheet number: 360

